\section{Discussion}
\label{sec:discussion}

\subsection{Why the results are what they are}
\label{sec:discussion-why}

Three mechanisms account for the measurements. The first is the spending of size. A fixed-window test spends its whole size at one reading; a monitor that is valid at every time spends it over all days, and by day 250 the rate e-process has spent about \nExactESize\,\% of a nominal 5\,\%. \Cref{prop:power} turns that into a bound on the power by day 250, and the measured power sits under the bound: \nByThreeE\,\% against a bound of \nExactThreeNpE\,\% at three times the nominal.
The power cost is therefore largely the consequence of the guarantee and not of a poor design; what design can change is how much of the bound is attained, and the single alternative at the right severity attains most of the bound that belongs to its own spending: it has spent \nExactSingleThreeSize\,\% of the size by day 250, more than the mixture, so the bound for a test of that size is \nExactSingleThreeNp\,\% at three times the nominal level, and it reaches \nAbByThreeSingleThree\,\%. The second mechanism is the mixture penalty of \cref{prop:growth}. A mixture of $K$ alternatives costs $\log K$ in the exponent of the e-process, which at $K=8$ is \nLogEight\ nats against a threshold of $\log 20=\nLogTwenty$ nats,
so the mixture needs about \nKlMixThree\ days where the matched bet needs \nKlSprtThree\ at three times the nominal level, and the ablation reproduces the order (\cref{tab:ablation-grid}). The third is what the bet knows. The published e-backtest of \citet{wang2022} starts with no bet and learns it from the past with a cap, which is consistent with its slower start (a bet that is estimated from few hits stays small): its median is \nMedThreeWwz\ days at three times the nominal and its power by day 250 is \nByThreeWwz\,\%. The uniform-prior Bayes factor spreads its weight over every hit rate in the unit interval, most of which are far from any hit rate that matters at a level of 1\,\%, and has a median of \nMedThreeBeta\ days.
Both are valid and conservative (\nNullSizeWwz\ and \nNullSizeBeta\,\% size), and both are slower than the mixture whose alternatives were placed where the departures of interest are.

The same mechanisms read the PD results. The product of $G$ grade mixtures pays the penalty $\log K$ once per grade, the average once plus $\log G$, while the product's growth rate is the sum of the grade growth rates and the average's is that of the best grade. At a small relative error the penalties dominate, which is consistent with the average detecting \nPdOneOneAvg\,\% against the product's \nPdOneOneProd\,\% at a factor of $1.1$, and at larger errors the growth dominates and the two reach \nPdOneTwoFiveProd\,\% at $1.25$. This is an interpretation through the bound of \cref{prop:growth}, not a separate theorem, and it applies to
the independent generator used here.

\subsection{Where the monitors do not help}
\label{sec:discussion-failure}

The monitors are slow when the departure is small, because the time to reject grows with the inverse of the divergence of the true hit rate from the null (\cref{tab:theory-delay}). At 1.5 times the nominal, which at a 99\,\% VaR is 1.5 exceedances per hundred days against one, more than half of the streams are not rejected in 1,000 days, and a fixed-window test read once is no better placed. The same holds for PD: a 5\,\% relative error in every grade is detected by the upper-side average in \nPdOneZeroFiveAvgUp\,\% of the ten-year histories, where the chi-square test at year 10 detects \nPdOneZeroFiveChiOne\,\%. A monitor that reads the PD of
a cohort every year is not a more sensitive instrument for a small bias than a test read once at the end; it is an instrument that warns early about a large one, and the data show that it does that (every history at $1.25$ by a median year of \nPdMedYearOneTwoFive). It is not more sensitive either than a chi-square rule read every year against a critical value set to the size the monitor measured (\nPdcMatchedSize\,\%), which rejects \nPdcOneZeroFiveChiMatched\,\% of the histories at $1.05$ and \nPdcOneOneChiMatched\,\% at $1.1$ (\nPdcOneZeroFiveChiSim\,\% and \nPdcOneOneChiSim\,\% at a nominal size of 5\,\%, \cref{tab:pd-equal-size}); what the monitor adds is a guarantee that holds at any year without a critical value calibrated for a horizon.

The monitors also fail where their null does. A correlated default history is not conditionally binomial, and \cref{prop:overdispersion} and the experiment of \cref{sec:results-pd} show the consequence: at an asset correlation of $0.03$ the two-sided e-processes reject a correct PD in \nRhoZeroThreeProd\ to \nRhoZeroThreeAvg\,\% of the histories. The same assumption underlies the chi-square test (\nRhoZeroThreeChiOne\,\% at year 10), so the finding is not that sequential monitoring is worse than a fixed test but that neither should be read against a binomial null for a portfolio whose defaults are correlated. A test whose null states the correlation keeps its size: the Vasicek-corrected binomial test of the engine rejects \nPdcZeroThreeVasTrueEvery\,\% of the histories at $0.03$ when it is given the true correlation, and loses it when it is given a correlation that is too low (\nPdcTwoFourVasAssumedEvery\,\% at $0.24$ with \nPdcAssumedRho\ assumed; \cref{tab:pd-comparators}). It keeps it with a correlation assumed above the truth only by having no power: with \nPdcAssumedRho\ assumed it rejects none of the histories of a PD that is $5\,\%$, $10\,\%$ or $25\,\%$ too low (\nPdcOneZeroFiveVasAssumedEvery, \nPdcOneOneVasAssumedEvery\ and \nPdcOneTwoFiveVasAssumedEvery\,\%; \cref{tab:pd-equal-size}), so the remedy holds only where the correlation is known, and the power of that test with the true correlation under a miscalibrated PD is not measured here. The sequential construction compounds the loss, because the process
multiplies one factor above its target per year, so a long history ends rejected while a short one may not be.

\subsection{What the results imply for a validation function}
\label{sec:discussion-use}

The results support three uses and argue against one. The monitors suit an early-warning role next to the supervisory tests. In the equity crisis they flag every model at a median of \nCElevenEquityNinetyEMed\ days where the traffic light needs 250, and in the foreign-exchange window they flag \nCElevenFxNinetyEFlag\ models at a median of \nCElevenFxNinetyEMed\ days, with an anytime-valid p-value at every step that a desk can read without inflating the size. They suit a calm period where the question is whether a model is being rejected by chance: in the control window the e-process flags \nCElevenCalmNinetyEFlag\ of 18 models where the traffic light read daily flags
\nCElevenCalmNinetyLightFlag\ and the Kupiec test on the last 250 days \nCElevenCalmNinetyKupTrailFlag. And they suit the reading of a vintage at every age, where the single reading at the end of life is too late for the cohorts that went wrong early (the 2006 to 2009 vintages are rejected at 12 months and their light at the mark is red).
The results argue against replacing the fixed-window tests: where a window is fixed in advance and read once, the likelihood-ratio test is the more powerful, and the monitor's own power cost at three times the nominal level is \nSeedsThreeCost\ points (\nSeedsUpThreeCost\ for the upper side).

For the choice of the configuration, the results give three rules. Use the upper-side monitor when the question is whether the VaR understates the risk, because it moves the delay and the power cost inside the criteria at both severities (a variant chosen after the two-sided results, so a recommendation that a prospective run has not yet confirmed; \cref{sec:discussion-negative}). Combine monitors of one stream by average, not by product. Set the grid of alternatives before the data to the range of departures that matter: a coarse grid is the cheapest improvement in power for the PD monitor, and
the default grid of the engine, fixed before these experiments, is the conservative choice.

\subsection{Negative and null results}
\label{sec:discussion-negative}

Five results go against the proposed monitors and stay in the paper. The delay of the two-sided monitor at three times the nominal level is \nSeedsThreeMed\ days, not below the window, and not below the traffic light's \nSeedsThreeLight\ either (the refutation clause of the design, on the real windows, does not fire). The power cost at that severity is \nSeedsThreeCost\ points, above the 10 allowed. The clustering monitor is weak against mild clustering: it rejects \nCluOneTenthCluster\,\% of the streams at $q=0.1$ and has a median beyond the horizon,
and the independence test on the last 250 days is faster at $q=0.4$ (a median of \nCluMedFourTenthIndep\ against \nCluMedFourTenthCluster\ days) at a size of \nNullSizeIndep\,\% when read daily. The PD monitor detects a 5\,\% relative error in \nPdOneZeroFiveAvg\,\% of the histories (the average) where a test read once detects \nPdOneZeroFiveChiOne\,\%, and its size is lost under correlation. At the size the monitor measured (\nPdcMatchedSize\,\%) the PD monitor also detects a small miscalibration less often than a chi-square rule does: for a PD $5\,\%$ too low the chi-square rule rejects \nPdcOneZeroFiveChiMatched\,\% of the histories, and \nPdcOneZeroFiveChiSim\,\% at a nominal size of 5\,\% (its size on these histories is \nPdcOneChiSim\,\%), against \nPdOneZeroFiveAvg\,\% for the monitor (\cref{tab:pd-equal-size}). A sixth observation concerns the upper-side variant, which fixes the first two failures; it was chosen after the two-sided results were known, so the pooled streams that support it test a configuration chosen by looking at the data, and a prospective run with fresh seeds is the confirmation it still needs.
